\chapter{Modernizing legacy locking and database access}
\label{ch:enterprise-migration}
\chapterbrief{Migrate locks while old and new writers coexist.}
\par\smallskip\noindent\textbf{Abbreviations used in this chapter:} DBI --- Database Interface; DDL --- data definition language.\par\smallskip

\index{Migration}
\index{MyISAM}
\index{Metadata locks}

Legacy applications often contain multiple generations of synchronization:
file locks around scripts, table locks around bulk writes, named advisory
locks around cross-table operations, and transactional row locks in newer
code. Removing a legacy lock without identifying every writer can create
a concurrency bug that did not previously exist.

\strategy{app-migration}{Incremental invariant migration}{Install the new guard before taking away the old one.}{multiple generations of writers share state.}{the contracts cannot coexist safely.}{DBI, InnoDB, regression tests}

\section{Inventory before migration}
\begin{figure}[!htbp]\centering
\begin{tikzpicture}[node distance=5mm]
\node[box,text width=10cm] (a) {Inventory writers, storage engines, locks and transaction owners};
\node[box,below=of a,text width=10cm] (b) {Measure contention and capture regression fixtures};
\node[okbox,below=of b,text width=10cm] (c) {Add guards, constraints and replay-safe identities};
\node[okbox,below=of c,text width=10cm] (d) {Canary new writers; compare and reconcile outcomes};
\node[box,below=of d,text width=10cm] (e) {Retire old locking only when every writer obeys the new contract};
\draw[arr](a)--(b); \draw[arr](b)--(c); \draw[arr](c)--(d); \draw[arr](d)--(e);
\end{tikzpicture}
\caption{Incremental migration preserves the invariant across mixed generations of code.}
\end{figure}
Inspect table engines, foreign keys, unique constraints, isolation levels,
queries, cron launchers and administrative scripts. MyISAM is nontransactional;
changing a DBI transaction boundary cannot make a MyISAM update roll back.
A move to InnoDB changes locking, index and operational requirements [68].
Test migration and rollback procedures on representative disposable data.

\section{Table locks and connection-local state}
MySQL \code{LOCK TABLES}, DDL and some administrative statements can
interact with transaction boundaries through implicit commits. Consult the
server's version-specific documentation [11,67]. Avoid combining
table-locking legacy code with ORM-managed transactions without a precise
contract. Named advisory locks, temporary tables, session variables and
isolation overrides are connection-local. A pooled connection must be reset
before reuse, or explicitly dedicated to the operation.

Long-running reports can retain snapshots and delay purge or contend with
writers, depending on the query and isolation. Prefer appropriate indexes,
bounded report windows and consistent snapshots. A replica may be useful
for reports that tolerate lag, but not for correctness-critical read-after-write
decisions unless the required consistency is enforced.

\warn{An online schema change can still require metadata locks. Test the
DDL algorithm and lock requirements for the exact server and operation;
do not assume that an online label means zero blocking.}

\section{Migration acceptance criteria}
Before removing a lock, list every code path that can mutate the protected
resource. Establish an invariant test, deploy the new guard while preserving
compatibility, and compare results with representative load. Record a
rollback strategy that does not reintroduce mixed incompatible writers.
Avoid a big-bang migration from local to distributed locking when a database
constraint can solve the problem more directly.

